\section{Conclusion}

We presented \textbf{Step-dLLM}, a sparse attention framework designed for the unique dynamics of diffusion language models. Our analysis reveals two properties that prior methods overlook: attention patterns in DLMs are locally consistent but not globally stationary across denoising steps, and different attention heads exhibit markedly different concentration behaviors. Step-dLLM exploits both observations through three integrated components: cross-step mask reuse with periodic refresh at anchor steps, a head-adaptive budget that allocates sparse computation according to each head's cumulative saliency, and fused Triton kernels that translate algorithmic savings into wall-clock speedups. Extensive experiments on LLaDA-1.5 and Dream-7B-Instruct across mathematical reasoning, code generation, and long-context retrieval benchmarks show that Step-dLLM consistently achieves the best accuracy--sparsity trade-off among sparse baselines, with advantages that grow under high sparsity---where stale masks are most harmful---and on long-context tasks, where the set of salient positions shifts substantially during denoising. At the kernel level, the fused block-sparse attention kernel delivers up to $8.6\times$ speedup over dense Flash Attention at 95\% sparsity.

Several directions remain open. The current anchor schedule uses a fixed refresh interval~$\Delta$; an adaptive strategy that triggers re-anchoring when the KL divergence between consecutive steps exceeds a threshold could better track abrupt attention transitions while further reducing the number of full-attention computations. Additionally, exploring whether the head-adaptive budget can be extended to full head pruning---entirely skipping heads whose saliency falls below a threshold at certain denoising stages---may yield further efficiency gains. Finally, integrating Step-dLLM with complementary acceleration techniques such as semi-autoregressive block decoding and step-reduction schedules is a promising avenue for compounding speedups in diffusion language model inference. More broadly, the anchor-refresh paradigm---periodically refreshing a sparse mask at selected steps to track shifting attention structure---is not specific to masked discrete diffusion: it may generalize to continuous diffusion transformers, multi-round editing, or any iterative generation process where attention patterns evolve across steps.
